\documentclass[../../../thesis.tex]{subfiles}
\begin{document}
\subsection{Výpisy kódov programu a algoritmy}
\label{sec:listings}
Ak je súčasť cieľov práce tvorba softvéru, prípadne analýza progрamátorských riešení,
je žiaduce uvádzať časti kódov vo forme krátkych výpisov (angl. \foreignlanguage{english}{\emph{listing}}).
Existuje niekoľko balíčkov, ktoré umožňujú zahrnúť časti kódov do textu práce.
Šablóna \verb|FEIstyle| používa balík \verb|listings|.
Prostredie \verb|lstlisting| vytvorí blok kódu s~menovkou Výpis kódu s~číslom výpisu.
Makro \verb|\FEIlistOfListings| na začiatku dokumentu vygeneruje zoznam všetkých výpisov,
ktorý nie je povinnou súčasťou práce,
býva však dobrým zvykom uvádzať ho najmä v~informatických študijných progрamoch.

Ukážka kódu je vo výpise~\ref{lst:main-c}.
V dodatku \ref{att:listings} nájdeme príklad výpisu obsahu externého textového súboru.
Ďalšie podrobnosti možno nájsť v~dokumentácii k~balíčkom%
\footnote{\href{https://ctan.org/pkg/listings}{\texttt{ctan.org/pkg/listings}}}
alebo v~tutoriáli služby Overleaf%
\footnote{\texttt{\href{https://www.overleaf.com/learn/latex/Code_listing}{www.overleaf.com/learn/latex/Code\_listing}}}.

\begin{lstlisting}[,
  caption={Ukážka výpisu kódu progрamu},
  label={lst:main-c},
  language=c,
  style=code-listing
]
/* Hello World program */

#include<stdio.h>

struct cpu_info {
    long unsigned utime, ntime, stime, itime;
    long unsigned iowtime, irqtime, sirqtime;
};

main()
{
    printf("Hello World");
}
\end{lstlisting}

\subsection*{Algoritmy}
Dvojica balíkov%
\footnote{\href{https://ctan.org/pkg/algorithms}{\texttt{ctan.org/pkg/algorithms}}}
\verb|algorithm|
a~\verb|algorithmic| uľahčuje zápis algoritmizácie pomocou tzv. pseudokódov.
Tieto nástroje sa uplatňujú najmä v prípade teoretických prác v~oblasti informatiky a~softvérového inžinierstva.
Šablóna FEIstyle načíta balíky automaticky a~tiež definuje makro \verb|\FEIlistOfAlgorithms|, ktoré na začiatku vytvorí zoznam všetkých algoritmov, v~záverečnej práci.
Príkaz možno vynechať alebo označiť riadok ako poznámku znakom \verb|%|, zoznam sa tak nevytvorí. Ako príklad uvádzame algoritmus \ref{alg:calculation-1} v~dodatku \ref{att:algorithms}.

Ďalšie podrobnosti získame z~tutoriálov%
\footnote{\href{https://www.overleaf.com/learn/latex/Algorithms}{\texttt{www.overleaf.com/learn/latex/Algorithms}}}
alebo z~dokumentácie k~jednotlivým balíčkom.
\end{document}